% ===========================================================================
% ΕΝΙΑΙΟ ΣΥΣΤΗΜΑ ΣΧΗΜΑΤΩΝ ΦΥΣΙΚΗΣ — styles/book-figures.tex
% ===========================================================================
% ΤΟ ΔΙΑΒΑΖΟΥΜΕ ΩΣ «ΛΕΞΙΚΟ» ΣΧΕΔΙΑΣΗΣ. Οι εντολές δεν δημιουργούν μόνες
% τους φυσικά σχήματα: ορίζουν εμφάνιση, όχι περιεχόμενο/γεωμετρία.
% Φορτώνεται από universalbook.sty ΜΕΤΑ την εκάστοτε παλέτα χρωμάτων.
% Αυτό σημαίνει ότι δεν γράφουμε ποτέ HTML/CMYK χρώματα στα chapters/*.tex.
%
% ΠΡΟΣΟΧΗ — ΕΚΚΡΕΜΕΙ ΕΓΚΡΙΣΗ ΣΧΗΜΑΤΩΝ.
% Οι υπάρχουσες μορφές είναι ΠΡΟΣΩΡΙΝΑ δείγματα από παλαιότερη εργασία,
% ΟΧΙ συμβάσεις που εγκρίθηκαν σε αυτό το project. Δεν έχει συμφωνηθεί
% χρωματικός κώδικας για ταχύτητα, επιτάχυνση, δυνάμεις, βάρος, τριβή,
% σχήματα/βέλη ή γενικός προσανατολισμός αξόνων.
% Κανένα AI δεν πρέπει να τις θεωρεί δεσμευτικές για νέο βιβλίο.
% Αρχή που έχει εγκριθεί: κοινή γεωμετρία/ίδια δεδομένα σε SCREEN και PRINT,
% PRINT δύο μελάνια. Οι υπόλοιπες προδιαγραφές μένουν ανοιχτές.
% Οι μορφοποιήσεις book/velocity, book/friction κ.λπ. παραμένουν μόνο
% για τεχνική συμβατότητα με τα ενδεικτικά demo μέχρι να τις ορίσει ο χρήστης.
% ---------------------------------------------------------------------------

% ---------------------------------------------------------------------------
% v0.27 — ΧΡΩΜΑΤΑ ΣΧΗΜΑΤΩΝ: ΜΟΝΟ ΣΗΜΑΣΙΟΛΟΓΙΚΟΙ ΡΟΛΟΙ
% Όλες οι τιμές BookFig*Color ορίζονται από το config/color-map.json
% και φορτώνονται μία φορά στο config/color-roles.tex.
% Έτσι το BookSecondary της ταχύτητας είναι το 55% και στο PRINT.
% ΑΠΑΓΟΡΕΥΟΝΤΑΙ εδώ κλάδοι SCREEN/PRINT και απευθείας ProofRed*.
% Τα ενδεικτικά χρώματα φυσικών μεγεθών ΔΕΝ είναι τελικά.
% Οι άξονες και η τριβή στα δείγματα χρησιμοποιούν προς το παρόν
% ουδέτερο BookInk σε SCREEN/PRINT (καμία κρυφή πορτοκαλί εξαίρεση).
% ---------------------------------------------------------------------------

% ---------------------------------------------------------------------------
% TikZ: χρήση \draw[book/axis] ...; ή \draw[book/velocity] ...;
% Οι διακεκομμένες γραμμές ερμηνεύονται στο υπόμνημα, όχι αποκλειστικά
% μέσω χρώματος. Οι αιχμές δεν αλλάζουν μεταξύ διαφορετικών συγγραφέων.
% ---------------------------------------------------------------------------
\tikzset{
  book/.cd,
  axis/.style={draw=BookFigAxisColor,->,line width=\BookFigAxisWidth,
               >={Stealth[length=\BookFigArrowLength,width=\BookFigArrowWidth]}},
  origin/.style={circle,fill=BookFigAxisColor,inner sep=1.3pt},
  body/.style={draw=BookFigObjectColor,fill=white,thick,line join=round},
  velocity/.style={draw=BookFigVelocityColor,->,line width=\BookFigVectorWidth,
                   >={Stealth[length=\BookFigArrowLength,width=\BookFigArrowWidth]}},
  accel/.style={draw=BookFigAccelColor,->,line width=\BookFigVectorWidth,dash pattern=on 6pt off 2pt,
                >={Stealth[length=\BookFigArrowLength,width=\BookFigArrowWidth]}},
  force/.style={draw=BookFigForceColor,->,line width=\BookFigVectorWidth,
                >={Stealth[length=\BookFigArrowLength,width=\BookFigArrowWidth]}},
  weight/.style={draw=BookFigWeightColor,->,line width=\BookFigVectorWidth,
                 >={Stealth[length=\BookFigArrowLength,width=\BookFigArrowWidth]}},
  friction/.style={draw=BookFigFrictionColor,->,line width=\BookFigVectorWidth,dash pattern=on 5pt off 1.8pt on 1pt off 1.8pt,
                   >={Stealth[length=\BookFigArrowLength,width=\BookFigArrowWidth]}},
  guide/.style={draw=BookFigGuideColor,line width=\BookFigAuxWidth,dashed},
  trajectory/.style={draw=BookFigVelocityColor,line width=\BookFigGraphWidth},
  plotA/.style={draw=BookFigPlotAColor,line width=\BookFigGraphWidth},
  plotB/.style={draw=BookFigPlotBColor,line width=\BookFigGraphWidth,dash pattern=on 6pt off 2.5pt},
  physlabel/.style={font=\small,inner sep=2pt,text=BookInk},
  endpoint/.style={circle,fill=BookFigObjectColor,inner sep=1.4pt}
}

% ---------------------------------------------------------------------------
% PGFPLOTS: κοινές διαστάσεις, μονάδες, άξονες, grid και διακριτές σειρές.
% \begin{axis}[book/plot, xlabel={$t\,(\mathrm{s})$},ylabel={$x\,(\mathrm{m})$}]
%   \addplot[book/seriesA] {1+0.5*x};
% \end{axis}
% Όταν το axis βρίσκεται μέσα σε tikzpicture, παραμένει ΔΙΑΝΥΣΜΑΤΙΚΟ.
% ---------------------------------------------------------------------------
\pgfplotsset{
  book/plot/.style={
    width=10.4cm,height=6.0cm,
    axis lines=left,
    axis line style={draw=BookFigAxisColor,line width=\BookFigAxisWidth,->},
    tick style={draw=BookFigAxisColor,line width=.5pt},
    ticklabel style={font=\scriptsize,color=BookInk},
    label style={font=\small,color=BookInk},
    xlabel style={at={(axis description cs:1.01,0)},anchor=west},
    ylabel style={at={(axis description cs:0,1.02)},anchor=south,rotate=0},
    grid=major,
    major grid style={draw=BookFigGuideColor,line width=.3pt,densely dotted},
    clip=false,
    enlarge x limits={upper,value=0.04},
    enlarge y limits={upper,value=0.06}
  },
  book/seriesA/.style={color=BookFigPlotAColor,very thick,solid,mark=none},
  book/seriesB/.style={color=BookFigPlotBColor,very thick,dashed,mark=none},
  book/seriesPoints/.style={only marks,mark=*,mark size=1.65pt,color=BookFigPlotAColor}
}

% ΚΕΝΤΡΙΚΟΙ ΚΑΝΟΝΕΣ ΓΙΑ ΦΩΤΟΤΥΠΙΑ / ΔΥΟ ΜΕΛΑΝΙΑ:
% - μη βασίζεσαι αποκλειστικά στο χρώμα όταν ξεχωρίζεις δύο φυσικά μεγέθη,
% - μην χρησιμοποιείς αυθαίρετα RGB μέσα σε έναν πίνακα ή διάγραμμα,
% - τα άκρα/μονάδες/βέλη να διαβάζονται χωρίς το χρωματικό υπόμνημα,
% - έλεγξε την απεικόνιση στην εκτύπωση 17x24 cm πριν την οριστικοποίηση.
